\section{Matched Simulation and Physical Architecture}
\label{sec:system}

\subsection{Design Principle: One Stack, Two Plants}

The experiment requires that any difference between simulated and physical
outcomes be attributable to the plant and its sensing, never to the autonomy
code. We therefore run a single implementation of the autonomy pipeline and
swap only what sits below it. Detection, temporal obstacle estimation, local
planning, and the body-velocity command interface are the same ROS~2 nodes
in both domains; the simulated plant is HoloOcean driving a native BlueROV2
Heavy agent, and the physical plant is a BlueROV2 Heavy with an ArduSub
autopilot. Fig.~\ref{fig:architecture} shows the two paths and the boundary
at which they diverge.

%==============================================================================
% Figure 1 --- matched sim/real architecture.  Status: READY.
% Source of truth: docs/SCIENTIFIC_BASELINE_0.md (architecture block),
% docs/PLANNER_INPUT_EQUIVALENCE.md, docs/SIM_REAL_SENSOR_EQUIVALENCE.md.
%==============================================================================
\begin{figure*}[t]
\centering
\begin{tikzpicture}[
  font=\scriptsize,
  node distance=2.6mm,
  box/.style={draw, rounded corners=1pt, align=center, minimum height=6.2mm,
              inner xsep=2pt, text width=17mm},
  shared/.style={box, fill=blue!8, very thick},
  simb/.style={box, fill=gray!7},
  realb/.style={box, fill=orange!8},
  valb/.style={box, fill=green!7, densely dashed},
  ar/.style={-{Latex[length=1.6mm]}, thick},
  arv/.style={-{Latex[length=1.6mm]}, densely dashed, gray!70!black},
]

% ---------------- shared autonomy column (center) -------------------------
\node[shared] (det)  {Object\\detector};
\node[shared, right=of det] (est) {Temporal\\estimator\\T0--T3};
\node[shared, right=of est] (pl)  {Local planner\\\PC{} or \PD{}};
\node[shared, right=of pl]  (tw)  {\texttt{cmd\_vel\_safe}\\(Twist)};

\node[above=1.2mm of est, font=\scriptsize\bfseries, text width=40mm,
      align=center] (sharedlab)
      {IDENTICAL ROS~2 CODE IN BOTH DOMAINS};

\draw[ar] (det) -- (est);
\draw[ar] (est) -- (pl);
\draw[ar] (pl)  -- (tw);

% ---------------- simulation row (top left / top right) -------------------
\node[simb, left=6mm of det, yshift=8.5mm] (simcam) {HoloOcean\\RGB camera};
\node[simb, left=of simcam] (simenv) {HoloOcean\\world +\\obstacle};
\node[simb, right=6mm of tw, yshift=8.5mm] (simctl)
      {Latency buffer,\\rate limit, PI\\vel.\ control};
\node[simb, right=of simctl] (simdyn)
      {8-thruster alloc.\\$\rightarrow$ BlueROV2\\rigid-body physics};

\draw[ar] (simenv) -- (simcam);
\draw[ar] (simcam) -- ++(0,-4.2mm) -| (det);
\draw[ar] (tw) -- ++(0,8.5mm) -- (simctl);
\draw[ar] (simctl) -- (simdyn);
\draw[ar] (simdyn.north) -- ++(0,3mm) -| node[pos=0.25, above, font=\tiny]
      {pose, attitude, depth (runtime subset)} (simenv.north);

% ---------------- real row (bottom left / bottom right) -------------------
\node[realb, left=6mm of det, yshift=-8.5mm] (realcam)
      {BlueROV2 camera\\H.264 1080p30};
\node[realb, left=of realcam] (realenv) {Pool +\\physical\\obstacle};
\node[realb, right=6mm of tw, yshift=-8.5mm] (realctl)
      {Safety interlock\\$\rightarrow$ ArduSub\\body-velocity};
\node[realb, right=of realctl] (realdyn)
      {8 thrusters\\$\rightarrow$ BlueROV2\\Heavy};

\draw[ar] (realenv) -- (realcam);
\draw[ar] (realcam) -- ++(0,4.2mm) -| (det);
\draw[ar] (tw) -- ++(0,-8.5mm) -- (realctl);
\draw[ar] (realctl) -- (realdyn);
\draw[ar] (realdyn.south) -- ++(0,-3mm) -| node[pos=0.88, below, font=\tiny]
      {\textsc{attitude}, pressure depth} (realenv.south);

% ---------------- validation-only lane ------------------------------------
\node[valb, below=12mm of pl, text width=34mm] (val)
      {\textbf{Validation only} --- never in the control path:\\
       simulator ground truth $\;\vert\;$ overhead RealSense};
\draw[arv] (simdyn.east) -- ++(3mm,0) |- (val.east);
\draw[arv] (realdyn.east) -- ++(3mm,0) |- (val.east);

\node[left=1mm of simenv, font=\scriptsize\bfseries, rotate=90,
      anchor=south] {SIM};
\node[left=1mm of realenv, font=\scriptsize\bfseries, rotate=90,
      anchor=south] {REAL};

\end{tikzpicture}
\caption{Matched architecture. The autonomy path (blue) is one
implementation executed unchanged in both domains; only the plant and its
sensing differ. Validation information (green, dashed) reaches the offline
validator and never the control path --- a property enforced by a
forbidden-topic guard in the estimator and by a fixed planner subscription
set, both unit-tested. The simulation-to-reality mismatches that the
calibration ladder of Sec.~\ref{sec:calibration} addresses live exactly on
the two vertical boundaries: camera and timing on the left, vehicle response
and sensing on the right.}
\label{fig:architecture}
\end{figure*}


\subsection{Runtime Information versus Validation Information}

A predictivity study is only meaningful if the simulated stack cannot see
anything its physical counterpart lacks, so the separation between runtime
and validation information is enforced rather than assumed. Runtime topics
carry quantities with a real counterpart: measured attitude, pressure depth,
the commanded body velocity, the image-space obstacle stream, and the
estimated odometry described in Sec.~\ref{sec:navigation}. Validation topics
--- simulator ground-truth pose, world-frame obstacle geometry, per-tick
dynamics traces in simulation, and the external overhead camera in the pool
--- are consumed only by the offline validator. The estimator node rejects
ground-truth-like input topics at construction, the planner subscription set
is fixed to three runtime topics, and unit tests assert both properties, so
a violation fails the build rather than silently improving a result.

\subsection{Sensing Equivalence and the Missing Velocity Sensor}

Matching the two domains required first auditing where they genuinely differ,
which surfaced one mismatch severe enough to change the system design. A
read-only inventory of the physical vehicle found no Doppler velocity log and
no velocity or position aiding of any kind in its telemetry: the available
runtime signals are attitude and rates, heading, external pressure depth, and
raw inertial measurements. The simulator, by contrast, offers a velocity
sensor and exact pose. Table~\ref{tab:sensors} records the resulting
equivalence audit. Two decisions follow. First, the simulator's velocity
sensor is confined to the simulated low-level controller, where it stands in
for the physical vehicle's own thruster-loop behavior, and is forbidden to
the navigation estimator. Second, runtime navigation is rebuilt around
signals that exist on both sides (Sec.~\ref{sec:navigation}). The remaining
entries in Table~\ref{tab:sensors} --- camera resolution and field of view,
transport latency, depth and attitude noise and rate --- are precisely the
quantities the calibration ladder of Sec.~\ref{sec:calibration} closes one
level at a time.

%==============================================================================
% Sim/real sensing equivalence.  Source: docs/SIM_REAL_SENSOR_EQUIVALENCE.md,
% docs/REAL_BLUEROV2_HARDWARE_INVENTORY.md.  Status: READY (real side is a
% completed read-only inventory; the "closed at" column is a plan, not a
% result).
%==============================================================================
\begin{table}[t]
\centering
\caption{Sensing equivalence audit between the simulated and physical
BlueROV2 Heavy, and the calibration level at which each residual mismatch is
closed. The absence of velocity aiding on the physical vehicle is a design
constraint, not a calibration target.}
\label{tab:sensors}
\footnotesize
\setlength{\tabcolsep}{3pt}
\renewcommand{\arraystretch}{1.2}
\begin{tabular}{@{}p{16mm} p{20mm} p{20mm} p{9mm}@{}}
\toprule
\textbf{Quantity} & \textbf{Simulation} & \textbf{Physical vehicle} &
\textbf{Closed at} \\
\midrule
RGB image      & HoloOcean camera, configurable resolution/FOV & H.264 USB camera, $1920{\times}1080$@30\,fps over UDP & \SONE \\
Image timing   & Tick-synchronous & Encode + transport + decode latency & \STWO \\
Attitude       & Pose sensor (noise-free) & \textsc{attitude} at 10\,Hz (EKF, IMU + compass) & \STHREE \\
Depth          & Depth sensor (noise-free) & External pressure sensor & \STHREE \\
Body velocity  & Velocity sensor (simulated plant only) & \textbf{none --- no DVL} & n/a \\
Position $x,y$ & Pose sensor (validation only) & \textbf{none at runtime} & n/a \\
Vehicle response & PI velocity loop + 8-thruster allocation & ArduSub body-velocity control, 8 thrusters & \STHREE \\
Contact        & Rigid-body collision & Physical contact & n/a \\
External GT    & Simulator pose (validation) & Overhead RealSense D435 (validation) & n/a \\
\bottomrule
\end{tabular}
\end{table}


\subsection{Simulation Interface and Timing}

Because the study measures latency and response, the simulator must advance
physics at the rate it claims to. Two timing defects were found and corrected
during integration and are reported here because they invalidate wall-clock
estimators silently. First, the host operating system's sleep quantization
capped the simulator serve loop below its nominal tick rate, so simulated
time ran slower than wall-clock time while ROS-side estimators kept
integrating; the loop now paces with a high-resolution timer and a
sleep/spin hybrid. Second, the underlying game engine advances physics by at
most $1/30$\,s per frame, so a nominal 20\,Hz tick rate advanced physics only
two thirds as far as declared and the vehicle moved at $0.646\times$ its
reported velocity. Scientific scenarios consequently run at exactly 30
ticks/s with verified real-time pacing. After both fixes, dead-reckoning
error on a straight 12.2\,m run fell from 4.08\,m to 0.117\,m maximum, which
confirms the estimator had been correct and the simulation clock had not.
This is a precondition for the timing calibration level \STWO{} and, more
generally, for any latency-aware sim-to-real claim.

The bridge between the simulator and ROS~2 is a two-process localhost
connection rather than the in-process bridge distributed with the
simulator, because the simulator and the ROS~2 distribution used here require
different Python interpreters on the experiment platform. This is an
infrastructure choice with no scientific content; it is documented for
reproducibility and is not claimed as a contribution.
